% from ega3-7-fr.tex
\oldpage[III]{195}
En effet, pour tout $x'\in\operatorname{Spec}(A')$, on a, en posant
$x=f(x')$,
\[
  H_\bullet(P_\bullet\otimes_A\kappa(x'))
  =
  H_\bullet\bigl((P_\bullet\otimes_A\kappa(x))
    \otimes_{\kappa(x)}\kappa(x')\bigr)
  =
  H_\bullet(P_\bullet\otimes_A\kappa(x))
    \otimes_{\kappa(x)}\kappa(x'),
\]
puisque $\kappa(x')$ est plat sur $\kappa(x)$, d'où la relation (7.6.7.1).
\end{proof}

\begin{lemma}[7.6.8]
Si l'anneau $A$ est noethérien et le complexe $P_\bullet$ formé de $A$-modules
de type fini, la fonction $x\mapsto d_p^T(x)$ sur $\operatorname{Spec}(A)$ est
constructible.
\end{lemma}

\begin{proof}
Il faut prouver que pour toute partie fermée irréductible $Y$ de
$X=\operatorname{Spec}(A)$, il existe un ouvert non vide $U$ de $Y$ dans lequel
$d_p$ est constante ($0$, 9.2.2); comme
$Y=\operatorname{Spec}(A/\mathfrak a)$, où $\mathfrak a$ est un idéal de $A$ tel
que $A/\mathfrak a$ soit réduit, on peut, en vertu de (7.6.7), se borner au cas
où $Y=X$ et où $A$ est un anneau noethérien intègre; mais alors l'assertion
résulte de (7.4.5).
\end{proof}

\begin{theorem}[7.6.9]
Soient $A$ un anneau noethérien, $P_\bullet$ un complexe de $A$-modules plats de
type fini, $T_\bullet(M)=H_\bullet(P_\bullet\otimes_A M)$ le foncteur homologique
défini par $P_\bullet$; pour tout $x\in\operatorname{Spec}(A)$, soit
\[
  d_p(x)=\operatorname{rang}_{\kappa(x)}T_p(\kappa(x)).
\]
Alors :
\begin{enumerate}
  \item[(i)] La fonction $d_p$ est constructible et semi-continue
    supérieurement dans $\operatorname{Spec}(A)$.
  \item[(ii)] Si $T_p$ est exact, $d_p$ est continue (donc localement constante)
    dans $\operatorname{Spec}(A)$; la réciproque est vraie lorsque l'anneau $A$
    est réduit.
\end{enumerate}
\end{theorem}

\begin{proof}
\begin{enumerate}
  \item[(i)] La première assertion a été démontrée dans (7.6.8). Pour prouver
    la seconde, il suffit donc ($0$, 9.2.4) de montrer que si $x'\ne x$ est une
    générisation de $x$ dans $\operatorname{Spec}(A)$, on a
    $d_p(x')\leq d_p(x)$. Or, il existe alors un anneau de valuation discrète
    $B$ et un morphisme
    \[
      f:\operatorname{Spec}(B)\longrightarrow\operatorname{Spec}(A)
    \]
    tels que, si $a$ désigne le point fermé de $\operatorname{Spec}(B)$ et $b$
    son point générique, on ait $f(a)=x$ et $f(b)=y$ (II, 7.1.9). En vertu de
    la formule (7.6.7.1), on voit qu'on est ramené à démontrer l'inégalité
    $d_p(a)\geq d_p(b)$ dans $\operatorname{Spec}(B)$; mais cela n'est autre que
    l'inégalité (7.6.5.1)\footnote{Le principe de démonstration de (i) par
    réduction au cas d'un anneau de valuation discrète nous a été communiqué
    oralement par Hironaka.}.
  \item[(ii)] La première assertion a déjà été démontrée (7.3.4). Pour démontrer
    la réciproque, utilisons le critère valuatif (7.6.2); compte tenu de la
    formule (7.6.7.1), on est donc ramené au cas où $A$ est un anneau de
    valuation discrète; mais comme $\operatorname{Spec}(A)$ ne comporte alors
    que deux points, l'hypothèse que $d_p$ est constante implique bien que $T_p$
    est exact, en vertu de (7.6.5).
\end{enumerate}
\end{proof}

\begin{corollary}[7.6.10]
Soient $A$ un anneau noethérien, $\mathfrak p_i$ $(1\leq i\leq r)$ ses idéaux
premiers minimaux, $k_i$ le corps résiduel de $A_{\mathfrak p_i}$
$(1\leq i\leq r)$.
\begin{enumerate}
  \item[(i)] Pour tout $x\in\operatorname{Spec}(A)$, il existe un indice $i$
    tel que
    \[
      d_p(x)\geq\operatorname{rang}_{k_i}T_p(k_i).
      \tag{7.6.10.1}\label{III.7.6.10.1.fr}
    \]
\end{enumerate}
En particulier, si $A$ est intègre et si $K$ est son corps des fractions, on a
\[
  d_p(x)\geq\operatorname{rang}_K T_p(K)
  \tag{7.6.10.2}\label{III.7.6.10.2.fr}
\]
pour tout $x\in\operatorname{Spec}(A)$.
